% Author: Izaak Neutelings (May 2018)
% Inspiration: https://tex.stackexchange.com/questions/113900/draw-polarized-light
\documentclass[border=3pt,tikz]{standalone}
\usepackage{amsmath} % for \text
\usepackage{tikz}
\usepackage{tikz-3dplot}
\usepackage{physics}
\tikzset{>=latex} % for LaTeX arrow head
\usepackage{xcolor}
\usepackage{amsmath}
\colorlet{myblue}{black!40!blue}
\colorlet{myred}{black!40!red}
\colorlet{vcol}{green!50!black}
\colorlet{Ecol}{orange!90!black}
\colorlet{EVcol}{orange!80!black!60}
\colorlet{Bcol}{violet!90}
\tdplotsetmaincoords{70}{110}
\begin{document}



% Electromagnetic wave - colored
\begin{tikzpicture}[tdplot_main_coords, scale=1.7, line cap=round, line join=round,
                    axis/.style={black, thick,->},
                    vector/.style={>=stealth,->}]

  
  \draw[line width=4, black!30] (-1, {-cos(60)}, {sin(60)}) -- (1, {-cos(60)}, {sin(60)}) -- (1, {cos(60)}, {-sin(60)}) --  (-1, {cos(60)}, {-sin(60)}) -- cycle;

  \draw[vcol, -latex, very thick] (0, 0, 0) -- (0, {2*sin(60)}, {2*cos(60)})  node[below] {$\boldsymbol{m}$}; 

  \draw[Bcol, -latex, very thick] (0, {-cos(60)}, {sin(60)}) --++ (0, 0, 0.3) node[above] {$\boldsymbol{B}$};
  \draw[Ecol, -latex, very thick] (0, {-cos(60)}, {sin(60)}) --++ (0, -0.5, 0) node[above] {$\boldsymbol{F}_1$};

  \draw[myred, -latex, thick] (1.2, {-0.7*cos(60)}, {0.7*sin(60)})  -- (1.2, {-0.3*cos(60)}, {0.3*sin(60)}) node [midway, xshift=-5] {$I$};
  
  \draw[Ecol, -latex, very thick] (1, 0, 0) --++ (0.5, 0, 0) node[below] {$\boldsymbol{F}_3$};
  \draw[Ecol, -latex, very thick] (-1, 0, 0) --++ (-0.5, 0, 0) node[above] {$\boldsymbol{F}_4$};

  \draw[Bcol, -latex, very thick] (0, {cos(60)}, {-sin(60)}) --++ (0, 0, 0.3) node[above] {$\boldsymbol{B}$};
  \draw[Ecol, -latex, very thick] (0, {cos(60)}, {-sin(60)}) --++ (0, 0.5, 0) node[below] {$\boldsymbol{F}_2$};
  \tdplotsetrotatedcoords{0}{90}{0}
  \tdplotdrawarc[tdplot_rotated_coords, thick, red, ->]{(0,0,0)}{1.3}{180}{120}{anchor=south west}{$\theta$};
  
  \draw[black] (1.2, {-cos(60)}, {sin(60)}) -- (4, {-cos(60)}, {sin(60)});
  \draw[black] (1.2, {cos(60)}, {-sin(60)}) -- (4, {cos(60)}, {-sin(60)});
  \draw[black, latex-latex] (3, {-cos(60)}, {sin(60)}) -- (3, {cos(60)}, {-sin(60)}) node [midway, left] {$a$};
  \draw[black] (1, {1.2*cos(60)}, {-1.2*sin(60)}) -- (1, {2*cos(60)}, {-2*sin(60)});
  \draw[black] (-1, {1.2*cos(60)}, {-1.2*sin(60)}) -- (-1, {2*cos(60)}, {-2*sin(60)});
  
  \draw[black, latex-latex](1, {1.7*cos(60)}, {-1.7*sin(60)}) -- (-1, {1.7*cos(60)}, {-1.7*sin(60)}) node [midway, below, xshift=5, yshift=2] {$b$};
  
  % MAIN AXES
  \draw[axis] (-2,0,0) -- (2.5,0,0) node[right] {$x$};
  \draw[axis] (0,-2,0) -- (0,2,0) node[right] {$y$};
  \draw[axis] (0,0,-2) -- (0,0,2) node[left] {$z$};
  

  % % draw (anti-)nodes
  % \foreach \iNode [evaluate={\iOffset=\iNode-1;}] in {1,...,\nNodes}{
  %   \ifodd\iNode \drawBNode \drawENode % E overlaps B
  %   \else        \drawENode \drawBNode % B overlaps E
  %   \fi
  % }

\end{tikzpicture}
\end{document}